% Generated draft: requires mathematical and source review.
\documentclass[11pt]{article}
\usepackage[margin=25mm]{geometry}
\usepackage{fontspec}
\setmainfont{Latin Modern Roman}
\usepackage{amsmath,amssymb,mathtools,unicode-math}
\setmathfont{Latin Modern Math}
\usepackage{enumitem,needspace,xurl,hyperref}
\setlength{\parindent}{0pt}
\setlength{\parskip}{5pt}
\setlength{\emergencystretch}{3em}
\newcommand{\N}{\mathbb{N}}
\newcommand{\Z}{\mathbb{Z}}
\newcommand{\Q}{\mathbb{Q}}
\newcommand{\R}{\mathbb{R}}
\newcommand{\C}{\mathbb{C}}
\newcommand{\F}{\mathbb{F}}
\newcommand{\A}{\mathbb{A}}
\newcommand{\PP}{\mathbb{P}}
\newcommand{\E}{\mathbb{E}}
\newcommand{\eps}{\varepsilon}
\newcommand{\dd}{\,\mathrm{d}}
\begin{document}
\section*{UnsolvedMath ID: 30005961}
\textbf{Source problem code:} catalogue entry
\subsection{Definitions and mathematical statement}
Let \(X\) be a smooth projective variety over \(\C\) with \(\dim X=3\), and let \(f\in\operatorname{Aut}(X)\) be a biregular automorphism. The supplied excerpt does not include the promised definitions of ``strict Calabi--Yau'' and ``primitive.'' Consequently, the following explicit interpretation is provisional and requires comparison with the complete source.

Under the simply connected convention, a strict Calabi--Yau threefold satisfies
\[
K_X\simeq\mathcal O_X,
\qquad \pi_1(X^{\mathrm{an}})=\{1\}.
\]
Here \(K_X\) denotes the canonical line bundle and \(X^{\mathrm{an}}\) the associated complex analytic manifold. The intended convention for ``strict'' is unresolved in the supplied excerpt.

Using the equivariance framework stated in the excerpt, call \(f\) primitive if there do not exist a normal projective variety \(B\) over \(\C\), a dominant rational map \(\pi\colon X\dashrightarrow B\), and a birational self-map \(f_B\in\operatorname{Bir}(B)\) such that
\[
0<\dim B<3,
\qquad \pi\circ f=f_B\circ\pi
\]
as rational maps. Whether the source additionally requires connected general fibers in this definition must be checked.

For a very ample Cartier divisor \(H\) on \(X\), define, for \(0\leq p\leq3\),
\[
d_p(f)=\lim_{n\to\infty}
\left(\bigl((f^n)^*(H^p)\cdot H^{3-p}\bigr)_X\right)^{1/n}.
\]
The excerpt states that these limits exist and are independent of \(H\). For the smooth projective, biregular setting, its entropy characterization is
\[
h_{\mathrm{top}}(f)=\max_{0\leq p\leq3}\log d_p(f),
\qquad h_{\mathrm{top}}(f)>0\ \Longleftrightarrow\ d_1(f)>1.
\]
Here \(h_{\mathrm{top}}\) is the topological entropy on \(X^{\mathrm{an}}\).

\textbf{Construction problem, conditional on the terminology above.} Construct further pairs \((X,f)\) consisting of a smooth strict Calabi--Yau threefold over \(\C\) and a primitive biregular automorphism satisfying \(d_1(f)>1\). The original request is to find ``many examples.'' Neither a required number of examples nor a criterion for distinctness or novelty is supplied. Thus the source gives an open-ended construction request, rather than a specified assertion about infinitely many isomorphism classes or a classification. The report describes its discussion as speculative rather than a collection of definite results; its broader standing framework allows normal projective varieties of dimension \(d\) and birational self-maps, whereas this target concerns smooth threefolds and biregular automorphisms.
\subsection{Short English statement}
Construct additional smooth strict Calabi--Yau threefolds over \(\C\) with positive-entropy automorphisms that admit no equivariant rational map to a positive-dimensional lower-dimensional variety. The precise conventions for strictness and primitivity, and the meaning of ``many'' or ``further,'' require review. The supplied curation labels the problem partially solved; this draft does not independently verify that assessment.
\subsection{Sources}
\begin{itemize}
\item Local UnsolvedMath registry, supplied complete record: permanent ID \(30005961\), exact source problem code \texttt{catalogue entry}; title: \emph{Positive-Entropy Automorphisms of Strict Calabi--Yau Threefolds}.
\item Supplied report citation: \emph{Algebraic Geometry: Wall Crossing and Moduli Spaces, Varieties and Derived Categories} (2025), Oberwolfach Reports, pp.~1787--1854, DOI \texttt{10.4171/owr/2024/32}. Original OWR report: \url{https://doi.org/10.4171/owr/2024/32}. The registry separately records proposed year 2024; the year discrepancy is unresolved here.
\item The excerpt identifies the contribution as Keiji Oguiso's \emph{Automorphisms of Calabi--Yau threefolds from algebraic dynamics and the second Chern class}. It presents Question~1 as a question after Dinh, Sibony, and D.-Q. Zhang, citing source references [6], [7], and [8] for background, variants, and known results. Those references are not bibliographically identified in the supplied record.
\item Supplied literature source, labeled ``Oguiso--Truong primary paper'': \url{https://arxiv.org/abs/1306.1590}. No further bibliographic details are supplied.
\item The excerpt attributes well-definedness of dynamical degrees to Dinh--Sibony [3], relative dynamical degrees to Dinh--Nguy\^en [2], with an additional reference [13], and the entropy formula to Gromov--Yomdin. Full entries for the numbered references are absent from the supplied record.
\item Literature assessment date: \texttt{2026-08-22}. This is supplied dated curation, not independently verified. Its status is ``partially solved'' and its assessment states that primitive positive-entropy automorphisms have been explicitly constructed on some strict Calabi--Yau threefolds, while the request for many examples is not a classification theorem. The supplied background explicitly requests independent verification of time-sensitive status and cited claims.
\item The registry labels the statement \texttt{corrected\_verified} and says it was reconstructed and checked against the official report. These are supplied provenance assertions, not verification performed for this draft.
\end{itemize}
\end{document}
